\documentclass[12pt,a4paper]{article}

\usepackage[czech]{babel}
\usepackage[T1]{fontenc}
\usepackage[utf8]{inputenc}

\usepackage{geometry}
\usepackage[most]{tcolorbox}

\geometry{margin=2.5cm}

\title{\textbf{Rozbor díla „Farma zvířat“}}
\author{}
\date{}

\begin{document}
\maketitle

%------------------------------------------------
\section*{Autor}

\begin{tcolorbox}[breakable]

\textbf{Autor:}
\begin{itemize}
\item George Orwell (vl. jménem Eric Arthur Blair)
\end{itemize}

\textbf{Název díla:}
\begin{itemize}
\item Farma zvířat
\end{itemize}

\textbf{Literární směr:}
\begin{itemize}
\item Světová literatura 1. poloviny 20. století, antiutopie, politická fikce, satirická próza
\end{itemize}

\textbf{Století tvorby:}
\begin{itemize}
\item 20. století
\end{itemize}

\textbf{Další autoři stejného směru + dílo:}
\begin{itemize}
\item Jaroslav Hašek – Osudy dobrého vojáka Švejka za světové války
\item Karel Čapek – Továrna na absolutno
\end{itemize}

\textbf{Další autorova díla:}
\begin{itemize}
\item 1984
\item Hold Katalánsku
\end{itemize}

\textbf{Specifický styl autora:}
\begin{itemize}
\item Jednoduchý a srozumitelný jazyk, důraz na jasné vyjádření myšlenky, realistická atmosféra a kritika totalitních režimů
\end{itemize}

\end{tcolorbox}

%------------------------------------------------
\section*{Charakteristika uměleckého textu}

\begin{tcolorbox}[breakable]

\textbf{Literární druh, žánr, forma:}
\begin{itemize}
\item Epika, próza, alegorická novela (politická satira)
\end{itemize}

\textbf{Typ vypravěče:}
\begin{itemize}
\item Er-forma, extradiegetický, heterodiegetický vypravěč
\end{itemize}

\textbf{Dominantní slohový postup:}
\begin{itemize}
\item Vyprávěcí
\end{itemize}

\textbf{Vysvětlení názvu díla:}
\begin{itemize}
\item Děj knihy se odehrává na farmě, kde hospodářská zvířata po revoluci převezmou vládu nad statkem.
\end{itemize}

\textbf{Aktuálnost díla:}
\begin{itemize}
\item Dílo je dodnes aktuální, protože ukazuje mechanismy fungování totalitních režimů, manipulaci a zneužití moci.
\end{itemize}

\textbf{Místo a čas děje:}
\begin{itemize}
\item Farma pana Jonese v Anglii. Čas není přesně určen, ale symbolicky odpovídá období vývoje Sovětského svazu (přibližně 1917–1943).
\end{itemize}

\textbf{Smysl díla:}
\begin{itemize}
\item Dílo ukazuje fungování totalitního režimu a postupnou proměnu revolučních ideálů v diktaturu. Zároveň varuje před propagandou, manipulací a zneužitím moci.
\end{itemize}

\textbf{Pravděpodobný adresát:}
\begin{itemize}
\item Široká veřejnost, především čtenáři v Evropě
\end{itemize}

\textbf{Kompozice:}
\begin{itemize}
\item Chronologická
\end{itemize}

\textbf{Zařazení do kontextu autorovy tvorby:}
\begin{itemize}
\item Dílo tematicky souvisí s románem 1984. Obě knihy kritizují totalitní režimy a manipulaci společnosti.
\end{itemize}

\textbf{Film nebo dramatizace:}
\begin{itemize}
\item Animal Farm (1954)
\item Film poměrně věrně zachycuje děj knihy a její hlavní myšlenky.
\end{itemize}

\end{tcolorbox}

%------------------------------------------------
\section*{Úryvek k rozboru}

Zvířata stála tiše ve stodole a poslouchala starého kance. Jeho hlas byl pomalý, ale silný. Mluvil o tom, jak lidé využívají práci zvířat a berou jim vše, co vyprodukují. „Člověk je náš skutečný nepřítel,“ řekl. „Když odstraníme člověka, budeme pracovat jen pro sebe.“  
Zvířata cítila směs naděje a neklidu. Některá si představovala svobodu, jiná se bála, co přijde potom.

\begin{tcolorbox}[breakable]

\textbf{Atmosféra úryvku:}
\begin{itemize}
\item Napjatá, zároveň plná naděje
\end{itemize}

\textbf{Postavy v úryvku:}
\begin{itemize}
\item Starý Major
\end{itemize}

\textbf{Charakteristika postav:}
\begin{itemize}
\item Níže
\end{itemize}

\textbf{Hlavní postava}
\begin{itemize}
\item Není jednoznačně určena
\end{itemize}

\textbf{Vedlejší postavy}
\begin{itemize}
\item Napoleon (prase): Inteligentní a mocichtivý kanec, který se stane diktátorem. Představuje Stalina.
\item Kuliš (prase): Čestný a schopný vůdce, který navrhuje stavbu větrného mlýna, ale je Napoleonem vyhnán. Představuje Trockého.
\item Pištík (prase): Mistr propagandy, který obhajuje Napoleonova rozhodnutí.
\item Starý Major (prase): Kanec, který inspiruje zvířata k revoluci. Symbolizuje myšlenky Marxe a Lenina.
\item Boxer (kůň): Pracovitý a oddaný kůň, který představuje pracující lid.
\item Benjamin (osel): Cynický a skeptický pozorovatel, který vidí realitu, ale nezasahuje.
\item Lupina (klisna): Mateřská postava, která se stará o ostatní zvířata.
\item Molina (klisna): Marnivá klisna, která dává přednost pohodlí před svobodou.
\item Psi: Napoleonova tajná policie.
\item Ovce: Dav, který slepě opakuje hesla.
\item Pan Jones: Původní majitel farmy, symbol svrženého cara.
\item Mojžíš (havran): Vypráví o „Cukrkandlové hoře“, symbol náboženství.
\end{itemize}

\textbf{Další postavy díla:}
\begin{itemize}
\item --
\end{itemize}

\textbf{Vztahy mezi postavami:}
\begin{itemize}
\item Zvířata naslouchají proslovu starého Majora a respektují jeho myšlenky.
\end{itemize}

\textbf{Zařazení úryvku do děje:}
\begin{itemize}
\item Začátek knihy – Starý Major pronáší svůj revoluční projev.
\end{itemize}

\textbf{Jazykové prostředky:}
\begin{itemize}
\item Přímá řeč
\item Symbolika
\item Alegorie
\end{itemize}

\end{tcolorbox}

%------------------------------------------------
\section*{Děj}

\subsection*{Část první: Revoluce na farmě}

Děj začíná na farmě pana Jonese, kde žijí různá hospodářská zvířata. Jednoho večera svolá starý kanec jménem \textbf{Starý Major} všechna zvířata do stodoly. Pronese projev, ve kterém vysvětluje, že lidé zvířata vykořisťují a berou jim výsledky jejich práce. Vyzývá zvířata, aby se jednou vzbouřila proti lidem a vytvořila společnost, kde si budou všechna zvířata rovna.  

Krátce poté Starý Major umírá, ale jeho myšlenky zůstávají mezi zvířaty. Revoluční ideje začnou rozvíjet především prasata \textbf{Napoleon} a \textbf{Kuliš}. Když farmář jednou zapomene zvířata nakrmit, zvířata se vzbouří a vyženou pana Jonese z farmy. Statek přejmenují na \textbf{Farmu zvířat}.  

Prasata vytvoří ideologii zvanou \textbf{animalismus} a na zeď stodoly napíší \textbf{sedm přikázání}, která mají zaručit rovnost všech zvířat.

\subsection*{Část druhá: Boj o moc}

Po revoluci zvířata pracují s velkým nadšením. Nejpracovitější je kůň \textbf{Boxer}, který opakuje heslo „Budu pracovat ještě víc“.  

Mezi prasaty Napoleonem a Kulišem vzniká spor o vedení farmy. Kuliš navrhuje postavit větrný mlýn, který by vyráběl elektřinu. Napoleon je proti, ale nakonec nechá Kuliše vyhnat pomocí psů, které si tajně vychoval.  

Napoleon se stává jediným vládcem farmy a postupně zavádí diktaturu.

\subsection*{Část třetí: Proměna režimu}

Postavení zvířat se postupně zhoršuje. Prasata si žijí v přepychu, zatímco ostatní zvířata musí stále více pracovat.  

Sedm přikázání se postupně mění tak, aby ospravedlnila chování prasat. Boxer, který celý život tvrdě pracoval, je po vyčerpání prodán do masokombinátu.  

Na konci příběhu začnou prasata chodit po dvou nohách a chovat se jako lidé. Poslední přikázání zní:  

„Všechna zvířata jsou si rovna, ale některá jsou si rovnější než jiná.“  

Ostatní zvířata už nedokážou rozeznat rozdíl mezi prasaty a lidmi.

%------------------------------------------------
\section*{Názor na dílo}

Dílo George Orwella na mě udělalo silný dojem. Ukazuje, jak se původně dobré myšlenky revoluce mohou postupně změnit v totalitní režim. Kniha velmi dobře vysvětluje principy propagandy, manipulace a zneužití moci.  

Jazyk knihy je jednoduchý a srozumitelný, přesto obsahuje hlubokou myšlenku. Dílo bych doporučil čtenářům, kteří se zajímají o politickou satiru, historii nebo společenské problémy.

\end{document}
